% Generated by report/make_restated.py from theory/A3_lkatyusha_lyapunov.tex: do not edit (edit the note and rerun).
\begin{restated}{Definition}{def:naive}{A3}{0}{naive accelerated SVRG}
The \emph{naive scheme} $\mathrm{N}(\theta_1,\eta,p)$ is \eqref{eq:K1}--\eqref{eq:K5} with $\theta_2=0$ and the step $\eta>0$ taken as a free parameter with $\eta\theta_1\le1$:
\begin{equation*}\tag{\ref*{eq:naive}}
\begin{aligned}
x_t&=\theta_1z_t+(1-\theta_1)y_t,\qquad g_t\text{ as in \eqref{eq:K2}},\qquad z_{t+1}=\frac{\eta\sigma x_t+z_t-\frac\eta Lg_t}{1+\eta\sigma},\\
y_{t+1}&=x_t+\theta_1(z_{t+1}-z_t),\qquad w_{t+1}=\begin{cases}y_t&\text{w.p. }p,\\ w_t&\text{w.p. }1-p.\end{cases}
\end{aligned}
\end{equation*}
We set $\nu:=\frac{\eta\theta_1}{1-\eta\theta_1}\in(0,\infty]$, with $\nu=\infty$ in the limiting case $\eta\theta_1=1$; this case is admissible only with exact gradients $g_t=\Fp(x_t)$ (Proposition~\ref{prop:agd-lc}), because the Young step of L3 (Lemma~\ref{lem:L3}) requires $\nu<\infty$ (Remark~\ref{rem:eta-theta2}). The potential is the same $\Psi_t$ as in \eqref{eq:lk-lyap} (with $W_t$ defined with an arbitrary constant $a_W\ge0$ in place of $\frac{\theta_2(1+\theta_1)}{p\theta_1}$, which is $0$ when $\theta_2=0$).
\end{restated}
